\section{Introduction}
\paragraph{For chemists: what an ML benchmark and a split are.}
A machine-learning (ML) benchmark for molecular properties is a fixed table of molecules with measured values (the \emph{labels}), a rule for dividing it into a training part and a held-out test part (the \emph{split}), and a metric. Models learn from the training molecules and are scored on how well they predict the held-out ones; leaderboards then rank models by this score. The split determines what the score means: a \emph{random} split puts similar molecules on both sides and rewards interpolation; a \emph{scaffold} split keeps molecules sharing a core ring system on one side only, testing generalisation to new chemotypes; a \emph{temporal} split trains on older and tests on newer data, mimicking prospective use. A \emph{noise ceiling} is the best score any model could reach given that the labels themselves are imprecise.

\paragraph{For ML readers: the chemistry of the labels.}
Most public potency data originate in ChEMBL~\cite{TODO-chembl}, a database manually curated from the literature. Potency is reported as $K_i$ (inhibition constant) or $\mathrm{IC}_{50}$ (concentration giving 50\% inhibition), both in molar units, and usually converted to the log-scale \emph{pChEMBL} value, $-\log_{10}$ of the molar value (a pChEMBL of 7 means 100\,nM). An \emph{assay} is one experimental protocol measuring a molecule against a target (enzyme, receptor or cell line) under specific conditions. A ChEMBL \emph{document} is one source publication (or deposited dataset) from which assays were extracted; different documents typically mean different laboratories, protocols and reagents. A learning \emph{task} here is a (target, measurement type) pair, e.g.\ human target X, $K_i$.

\paragraph{Problem.}
A molecule tested in two documents often receives different values~\cite{TODO-landrum-riniker}. If two models differ by less than this disagreement, a leaderboard gap says little about real predictive quality. Prior work estimates performance bounds by adding noise of an assumed size~\cite{TODO-fitting-data-or-noise}; we instead \emph{measure} the noise from actual replicate measurements across independent documents.

\paragraph{Contributions.}
(i) A replicate-based estimator of per-endpoint noise ceilings with molecule-level bootstrap and sensitivity analysis. (ii) A fraction-of-ceiling metric that makes scores comparable across endpoints. (iii) A dataset of replicated molecules decontaminated against existing benchmarks. (iv) A benchmark comparing tuned small models and frozen foundation-model embeddings under three splits and five seeds. (v) Recommendations: minimum detectable effect, a reporting protocol and a checklist. Whether any claimed advantage of pretrained embeddings survives is \tbd{}.

\paragraph{Hypotheses (to be tested, with refutation criteria).}
H1: between-document noise is comparable to model error on a substantial share of endpoints (refuted if ceilings are near perfect; then the noise argument does not hold and the paper is reframed as a model comparison). H2: the advantage of pretrained embeddings over ECFP+GBM shrinks or vanishes when measured as fraction of ceiling (refuted if it persists across splits with paired confidence intervals excluding zero). Both are untested.
